\hypertarget{pd15:evidence}{%
\chapter{Public evidence for a causal banking
proposal}\label{pd15:evidence}}

A bank should invest in a causal-analysis service only when there is a
credible account of the decisions it could improve, the methods worth
assembling and the experiments worth funding. Public evidence can establish
technical promise and identify consequential uncertainties before a bank
pays to collect private assignment histories or commission a larger trial.
The proposed studies follow that investment sequence. This chapter distinguishes the evidence classes and the original survey's coverage; the next compares dataset families, and the following chapter maps them to a bounded study portfolio. Later completed studies are linked explicitly so that a historical acquisition decision is not mistaken for the current state of the scientific evidence.

Public data can provide much of the evidence for that first decision. It
can expose realistic variable definitions, missingness, reporting
clocks, financial magnitudes and computational workloads. Public
experiments can exercise estimation under a documented assignment
design. Semi-synthetic data can combine empirical covariates with
generated treatment assignments and potential outcomes, making error
measurable. Entirely synthetic histories can isolate feedback,
confounding, graph equivalence and observation errors. These sources
complement one another because no single source supplies realistic bank
operations, observed counterfactual truth and unrestricted reuse at
once.

Public studies can establish technical merit and identify a worthwhile
next experiment without estimating the effect of a policy in the adopting
bank. An effect measured in another population is not the bank's effect.
A negative comparison, an informative finding of nonidentification, or a
narrower useful capability can still improve the investment decision.

\hypertarget{what-constitutes-evidence}{%
\section{What constitutes evidence}\label{what-constitutes-evidence}}

Let \(S=1\) denote membership in a public study and \(S=0\) membership
in the future bank population. A public randomized study may identify

\[
\tau_1=\mathbb E[Y(1)-Y(0)\mid S=1].
\]

It does not by itself identify \(\tau_0\). One possible transport
argument would require comparable interventions and outcomes, sufficient
effect modifiers \(X\), overlap between populations, and stability of
the conditional causal response. Under those additional assumptions,

\[
\tau_0=\int\{m_1(x)-m_0(x)\}\,dP(x\mid S=0),
\qquad m_a(x)=\mathbb E[Y\mid A=a,X=x,S=1].
\]

The expression explains the later data requirement: a target population
distribution and a credible account of effect modification. It is not a
finding that these assumptions hold. A banking example makes the
distinction concrete. Learning how advertising holdouts support an
incremental-conversion estimate can validate treatment coding and an
estimator's workflow. It does not show the deposit value of a bank
telephone campaign, whose product, population, contact policy and
financial payoff differ.

Five evidence classes are kept distinct throughout this survey.

\begin{longtable}[]{@{}>{\raggedright\arraybackslash}p{\dimexpr 0.22\linewidth-2\tabcolsep\relax}>{\raggedright\arraybackslash}p{\dimexpr 0.36\linewidth-2\tabcolsep\relax}>{\raggedright\arraybackslash}p{\dimexpr 0.42\linewidth-2\tabcolsep\relax}@{}}
\toprule
\begin{minipage}[b]{\linewidth}\raggedright
Evidence class\strut
\end{minipage} & \begin{minipage}[b]{\linewidth}\raggedright
What it can demonstrate\strut
\end{minipage} & \begin{minipage}[b]{\linewidth}\raggedright
What it cannot establish alone\strut
\end{minipage}\tabularnewline
\midrule
\endhead
\begin{minipage}[t]{\linewidth}\raggedright
Public financial observations\strut
\end{minipage} & \begin{minipage}[t]{\linewidth}\raggedright
Data semantics, empirical associations, support, sensitivity, realistic
workflows and resource use\strut
\end{minipage} & \begin{minipage}[t]{\linewidth}\raggedright
Random assignment, omitted-information absence or target-bank
effects\strut
\end{minipage}\tabularnewline
\begin{minipage}[t]{\linewidth}\raggedright
Public randomized or designed studies\strut
\end{minipage} & \begin{minipage}[t]{\linewidth}\raggedright
Effects in the documented study population under its design; estimator
comparisons with appropriate uncertainty\strut
\end{minipage} & \begin{minipage}[t]{\linewidth}\raggedright
Individual counterfactual labels, automatic transport or bank ROI\strut
\end{minipage}\tabularnewline
\begin{minipage}[t]{\linewidth}\raggedright
Semi-synthetic benchmarks\strut
\end{minipage} & \begin{minipage}[t]{\linewidth}\raggedright
Error against generated truth on realistic covariates, with the
generator disclosed\strut
\end{minipage} & \begin{minipage}[t]{\linewidth}\raggedright
Truth of the generated response model in a bank\strut
\end{minipage}\tabularnewline
\begin{minipage}[t]{\linewidth}\raggedright
Controlled synthetic histories\strut
\end{minipage} & \begin{minipage}[t]{\linewidth}\raggedright
Failure mechanisms, refusal behavior, exact references and
repeated-sampling calibration in a specified law\strut
\end{minipage} & \begin{minipage}[t]{\linewidth}\raggedright
Empirical prevalence of those mechanisms or authentic audit
accuracy\strut
\end{minipage}\tabularnewline
\begin{minipage}[t]{\linewidth}\raggedright
Public reasoning and linkage tasks\strut
\end{minipage} & \begin{minipage}[t]{\linewidth}\raggedright
Logical consistency, handling ambiguity, provenance and observation
construction\strut
\end{minipage} & \begin{minipage}[t]{\linewidth}\raggedright
Business ownership, real customer identities or live-model reliability
without actual evaluation\strut
\end{minipage}\tabularnewline
\bottomrule
\end{longtable}

For an observational file, a fitted model is not an assignment
mechanism. For a randomized experiment, the difference in sample means
is an estimate with uncertainty, not the true individual treatment
effect. For a semi-synthetic file, both potential outcomes may be
available, but they were generated. These distinctions determine which
metrics are meaningful: prediction accuracy for a repayment label,
interval coverage only when truth or a valid evaluation construction is
available, and effect identification only under explicit substantive
assumptions.

\hypertarget{survey-method-and-coverage}{%
\section{Survey method and coverage}\label{survey-method-and-coverage}}

The survey inspected 24 dataset families using custodian descriptions,
author-maintained benchmark documentation, data dictionaries and
selected loader/generator source. It covers funding and bank panels;
household finance and marketing; lending and mortgage histories;
experimental and semi-synthetic estimation; discovery; linkage; and
language-model reasoning. A source register preserves exact downloaded
bytes, SHA-256 identities, inspected passages, supported claims and
access limitations. Dataset documentation supports dataset descriptions;
this survey does not infer a scholarly result from a paper title or
repository abstract.

The publication cutoff is 7 October 2026. Eight benchmark documentation
files and the CFPB historical schema are pinned to repository commits at
or before that cutoff. Other web pages are acquisition snapshots from 12
October; their current text is not represented as an archived 7 October
page. Older survey and data releases described by those pages remain
candidates, but their exact selected files must be frozen before
analysis. The SCF page provides a concrete trap: some summary extracts
were updated on 9 October 2026, after the cutoff. Those extracts are
excluded from a cutoff-respecting study until a prior version is
recovered. The full 2022 dataset is described separately with an earlier
release date \citep{pd15_scf}.

This is a structured, bounded survey. It is strongest for US regulatory
and mortgage data, two commonly used non-US bank datasets, and public
methodological benchmarks. It is not an exhaustive international census.
China-specific supervisory and household panels, European credit
registers, financial inclusion and microcredit trial archives, and
detailed commercial panel providers need separate custodian and access
reviews. Their absence does not establish that no suitable public source
exists. No registration-gated mortgage account was created, no
commercial data agreement was accepted, and no linked external benchmark
code was executed during this survey.

\hypertarget{evidence-already-available}{%
\section{Evidence already available}\label{evidence-already-available}}

The repository already contains a reproducible public 401(k) analysis
and a quarterly monetary analysis, alongside its synthetic banking
studies. The 401(k) extract contains 9,915 households, eligibility,
participation, financial assets and nine declared controls. The existing
analysis distinguishes an eligibility average effect from a
participation instrumental-variable target and compares score algebra
with DoubleML under identical fitted inputs. The extract lacks
survey-design fields sufficient to establish design-valid covariance;
conditional exchangeability, instrument assumptions and bank transport
are unresolved. These are informative prototype limitations, not reasons
to discard the example \citep{public401k}.

The monetary example uses 203 quarterly observations from 1959Q1 through
2009Q3. It reproduces local-projection coefficients and Bartlett
covariance against a numerical reference, and makes the
information-clock ambiguity explicit. An aggregate money-stock response
to a Treasury-bill rate change is not a bank's deposit-offer effect. It
is nevertheless a useful public demonstration of time-series
conventions, dependence-aware computation and sensitivity to period and
control choices \citep{public_macrodata_dynamic}.

A new descriptive inspection downloaded the UCI Bank Marketing archive,
checked its four CSV files and preserved their identities. The two full
datasets contain 41,188 and 45,211 observations; their smaller files
contain 4,119 and 4,521. The 41,188-row version has 4,640 subscriptions
and 36,548 non-subscriptions. Its \texttt{default} field contains 8,597
\texttt{unknown} values and \texttt{education} contains 1,731.
Empty-cell counts alone would therefore conceal missing information. The
full and smaller versions overlap and must not be used as independent
replications \citep{pd15_uci_bank,pd15_profile}.

That inspection is a completed data-profile exercise. It is not a new
treatment-effect estimate, a fitted benchmark or a preregistered
experiment. It establishes that an openly licensed, modest-size
financial dataset is immediately usable for schema, timing and refusal
demonstrations. The subsequent public NSW analysis in Chapter~\ref{gap29:public} supplies a fitted experimental comparison under its stated selection and design premises. The panel, servicing and linkage studies in Chapters~\ref{gap29:funding}, \ref{gap29:servicing} and~\ref{gap29:observation} instead use generated laws. These additions extend the evidence without changing the original survey's acquisition record or supplying bank transport and audit truth.

